% frontmatter.tex -- Module 6 (Electromagnetism): cover, How to use this guide,
% Contents, Module overview.
% \inquiry boxes of ch1.tex ... ch5.tex (Chapters 3 and 4 share IQ3).
% "Formulas at a glance" = one row per formula card in ch1.tex ... ch5.tex (26 cards), in
% card order, with the card's own notation and badge (on the sheet / not on the sheet /
% Experiment list: the Key experiment boxes of ch1.tex ... ch5.tex.
%
% Local layout fix (common/ is shared, so it is made here, inside a group; same fix as
% Module 8): three-digit page numbers in the contents were 2 pt wider than the
% page-number box (titletoc \@pnumwidth = 1.55em), so they stuck out into the margin.
% The box is made 1.9em wide; the correction 0.65em keeps the text margin at 2.55em.
\makecover
\IfFileExists{../common/howtouse.tex}{% howtouse.tex -- shared "How to use this guide" pages (same in all four books).
% Input it from frontmatter.tex:   % howtouse.tex -- shared "How to use this guide" pages (same in all four books).
% Input it from frontmatter.tex:   % howtouse.tex -- shared "How to use this guide" pages (same in all four books).
% Input it from frontmatter.tex:   \input{../common/howtouse}
% Numbers in the mini worked example: R = 12 / 0.50 = 24 ohm (checked in Python).
\unnumberedchapter[\faInfoCircle]{How to use this guide}

This guide explains Module~\modulenumber\ (\moduletitle) from the very beginning.
It uses simple English, short sentences and many examples. Every calculation is
shown step by step.

The chapters follow the inquiry questions of the NSW Physics Stage 6 Syllabus
(2017). At the back of the book you will find a formula summary, a glossary
(with Chinese translations), a practice set, full answers and a last-minute
summary. In the PDF, you can click any entry in the contents, any
cross-reference and any question number.

\section*{What the coloured boxes mean}

Each kind of box has its own colour and icon. A colour always means the same
thing, in every book of this series.

\begin{sidebyside}[0.5]
  \begin{bigidea}[before skip=0pt, after skip=0pt, equal height group=legend1, fontupper=\normalsize]
    The main point of a topic in two or three sentences. Read these first.
  \end{bigidea}
\nextside
  \begin{keywords}[before skip=0pt, after skip=0pt, equal height group=legend1]
    \kw{Term}{what the word means, in plain words.}
    New words are also printed in \term{bold} the first time they appear.
  \end{keywords}
\end{sidebyside}

\begin{sidebyside}[0.5]
  \begin{mistakes}[before skip=0pt, after skip=0pt, equal height group=legend2]
    Errors that students often make, and how to avoid them:
    \xmark\ wrong \enspace \cmark\ right.
  \end{mistakes}
\nextside
  \begin{examtip}[before skip=0pt, after skip=0pt, equal height group=legend2]
    How the HSC asks about this topic, and what you must write to get full
    marks.
  \end{examtip}
\end{sidebyside}

\begin{sidebyside}[0.5]
  \begin{toolkit}[before skip=0pt, after skip=0pt, equal height group=legend3]{}
    A short maths skill that you need for the next part, for example
    rearranging a formula.
  \end{toolkit}
\nextside
  \begin{refresher}[before skip=0pt, after skip=0pt, equal height group=legend3]{}
    A short reminder of a Year~11 idea that you need, for example Newton's
    laws.
  \end{refresher}
\end{sidebyside}

\begin{sidebyside}[0.5]
  \begin{tcolorbox}[hsc minor=hscqc, unbreakable, before skip=0pt, after skip=0pt,
      equal height group=legend4, title={\hscboxtitle{\faQuestionCircle}{Quick check}{}}]
    Short questions to test yourself, numbered like \textcolor{hscqc}{\bfseries Q2.3}
    (Chapter~2, question~3). The answers at the back use the same numbers.
  \end{tcolorbox}
\nextside
  \begin{experiment}[unbreakable, before skip=0pt, after skip=0pt, equal height group=legend4]{}
    A practical or historical experiment: aim, apparatus, method, results,
    conclusion and why it matters.
  \end{experiment}
\end{sidebyside}

\begin{formulacard}[unbreakable, before skip=6pt]{Name of the formula}
  \cardformula{V = IR}
  Every formula has a card: the formula, the meaning and SI unit of each
  symbol, when to use it, and a common trap. The badge under the formula tells
  you about the official HSC formula sheet:
  \begin{itemize}
    \item \hscbadgeonsheet\ \ it is printed on the sheet, so you are given it
          in the exam;
    \item \hscbadgenotonsheet\ \ you must learn it, or know how to get it;
    \item \hscbadgecombine{two formulas}\ \ it is two formulas from the sheet
          joined together.
  \end{itemize}
\end{formulacard}

\begin{workedexample}[nonumber, unbreakable]{How every worked example is set out}
  A \qty{12}{V} battery drives a current of \qty{0.50}{A} through a resistor.
  Find the resistance. (Cover everything below the \emph{Solution} line with a
  sheet of paper and try it first.)
  \given{$V = \qty{12}{V}$, \ $I = \qty{0.50}{A}$}
  \find{the resistance $R$}
  \formula{$V = IR$, \ so $R = \dfrac{V}{I}$}
  \substitute{$R = \dfrac{\qty{12}{V}}{\qty{0.50}{A}}$ \ (numbers with their units)}
  \answer{\final{$R = \qty{24}{\ohm}$} \ (with a unit and sensible significant
    figures)}
  \checkstep{is the answer sensible? Units: $\unit{V} \div \unit{A} = \unit{\ohm}$ \cmark}
\end{workedexample}

Each chapter starts with its \textbf{inquiry question} from the syllabus, a
list of \textbf{goals} and the \textbf{syllabus dot points} it covers. It ends
with a short \textbf{summary}.

\section*{How to study with this guide in the last two weeks}

The HSC Physics exam is on \textbf{Thursday 5 November 2026}. You have four
books (Modules 5 to 8). Here is a simple plan.

\begin{table}
  \centering
  \begin{tabular}{@{}L{30mm}L{\dimexpr\linewidth-34mm\relax}@{}}
    \toprule
    \colhead{When} & \colhead{What to do}\\
    \midrule
    Days 1 to 10 & Work through the four books, about two or three days for
      each book. For every chapter, follow steps 1 to 5 below.\\
    Days 11 to 13 & Do the Practice set of each book like a real exam: time
      yourself and use only the formula sheet. Mark your answers with the
      marking criteria. Redo every question you got wrong.\\
    The night before & Read the Last-minute summary of each book and look at
      the formula sheet once more. Then go to bed early.\\
    \bottomrule
  \end{tabular}
\end{table}

For each chapter:
\begin{enumerate}
  \item Read the \textbf{big idea} and the \textbf{key words} (about 10 minutes).
  \item Read the explanation and the \textbf{formula cards}. Say each formula
        ``in words''.
  \item \textbf{Cover the solution} of each worked example and try it yourself.
        Then compare your work with the book, step by step.
  \item Do the \textbf{Quick check} questions. Mark them with the answers at the
        back.
  \item Write every mistake you made in a ``mistakes list''. Read this list
        again before the exam.
\end{enumerate}

\section*{How to use the official formula sheet in the exam}

In the exam you get the NESA data sheet and formulae sheet. It gives many
formulas and constants, but not all of them.

\begin{enumerate}
  \item \textbf{Before the exam}, learn where each formula is on the sheet. The
        sheet has four parts: motion, forces and gravity; waves and
        thermodynamics; electricity and magnetism; quantum, special relativity
        and nuclear. Then you can find a formula quickly.
  \item \textbf{Learn the formulas that are not on the sheet.} Their formula
        cards have the red badge and show you how to derive them.
  \item \textbf{Check what each symbol means} in the question. The same letter
        can mean different things: for example, $T$ is a period in one formula
        and a temperature in another.
  \item \textbf{Change every quantity to SI units first}: km to m, nm to m,
        \si{km.h^{-1}} to \si{m.s^{-1}}, eV to J.
  \item \textbf{Use the values on the data sheet}, for example
        $g = \qty{9.8}{m.s^{-2}}$ and $c = \qty{3.00e8}{m.s^{-1}}$.
  \item \textbf{Show your working} in the same order as the worked examples:
        formula, substitution, answer with a unit.
\end{enumerate}

% Numbers in the mini worked example: R = 12 / 0.50 = 24 ohm (checked in Python).
\unnumberedchapter[\faInfoCircle]{How to use this guide}

This guide explains Module~\modulenumber\ (\moduletitle) from the very beginning.
It uses simple English, short sentences and many examples. Every calculation is
shown step by step.

The chapters follow the inquiry questions of the NSW Physics Stage 6 Syllabus
(2017). At the back of the book you will find a formula summary, a glossary
(with Chinese translations), a practice set, full answers and a last-minute
summary. In the PDF, you can click any entry in the contents, any
cross-reference and any question number.

\section*{What the coloured boxes mean}

Each kind of box has its own colour and icon. A colour always means the same
thing, in every book of this series.

\begin{sidebyside}[0.5]
  \begin{bigidea}[before skip=0pt, after skip=0pt, equal height group=legend1, fontupper=\normalsize]
    The main point of a topic in two or three sentences. Read these first.
  \end{bigidea}
\nextside
  \begin{keywords}[before skip=0pt, after skip=0pt, equal height group=legend1]
    \kw{Term}{what the word means, in plain words.}
    New words are also printed in \term{bold} the first time they appear.
  \end{keywords}
\end{sidebyside}

\begin{sidebyside}[0.5]
  \begin{mistakes}[before skip=0pt, after skip=0pt, equal height group=legend2]
    Errors that students often make, and how to avoid them:
    \xmark\ wrong \enspace \cmark\ right.
  \end{mistakes}
\nextside
  \begin{examtip}[before skip=0pt, after skip=0pt, equal height group=legend2]
    How the HSC asks about this topic, and what you must write to get full
    marks.
  \end{examtip}
\end{sidebyside}

\begin{sidebyside}[0.5]
  \begin{toolkit}[before skip=0pt, after skip=0pt, equal height group=legend3]{}
    A short maths skill that you need for the next part, for example
    rearranging a formula.
  \end{toolkit}
\nextside
  \begin{refresher}[before skip=0pt, after skip=0pt, equal height group=legend3]{}
    A short reminder of a Year~11 idea that you need, for example Newton's
    laws.
  \end{refresher}
\end{sidebyside}

\begin{sidebyside}[0.5]
  \begin{tcolorbox}[hsc minor=hscqc, unbreakable, before skip=0pt, after skip=0pt,
      equal height group=legend4, title={\hscboxtitle{\faQuestionCircle}{Quick check}{}}]
    Short questions to test yourself, numbered like \textcolor{hscqc}{\bfseries Q2.3}
    (Chapter~2, question~3). The answers at the back use the same numbers.
  \end{tcolorbox}
\nextside
  \begin{experiment}[unbreakable, before skip=0pt, after skip=0pt, equal height group=legend4]{}
    A practical or historical experiment: aim, apparatus, method, results,
    conclusion and why it matters.
  \end{experiment}
\end{sidebyside}

\begin{formulacard}[unbreakable, before skip=6pt]{Name of the formula}
  \cardformula{V = IR}
  Every formula has a card: the formula, the meaning and SI unit of each
  symbol, when to use it, and a common trap. The badge under the formula tells
  you about the official HSC formula sheet:
  \begin{itemize}
    \item \hscbadgeonsheet\ \ it is printed on the sheet, so you are given it
          in the exam;
    \item \hscbadgenotonsheet\ \ you must learn it, or know how to get it;
    \item \hscbadgecombine{two formulas}\ \ it is two formulas from the sheet
          joined together.
  \end{itemize}
\end{formulacard}

\begin{workedexample}[nonumber, unbreakable]{How every worked example is set out}
  A \qty{12}{V} battery drives a current of \qty{0.50}{A} through a resistor.
  Find the resistance. (Cover everything below the \emph{Solution} line with a
  sheet of paper and try it first.)
  \given{$V = \qty{12}{V}$, \ $I = \qty{0.50}{A}$}
  \find{the resistance $R$}
  \formula{$V = IR$, \ so $R = \dfrac{V}{I}$}
  \substitute{$R = \dfrac{\qty{12}{V}}{\qty{0.50}{A}}$ \ (numbers with their units)}
  \answer{\final{$R = \qty{24}{\ohm}$} \ (with a unit and sensible significant
    figures)}
  \checkstep{is the answer sensible? Units: $\unit{V} \div \unit{A} = \unit{\ohm}$ \cmark}
\end{workedexample}

Each chapter starts with its \textbf{inquiry question} from the syllabus, a
list of \textbf{goals} and the \textbf{syllabus dot points} it covers. It ends
with a short \textbf{summary}.

\section*{How to study with this guide in the last two weeks}

The HSC Physics exam is on \textbf{Thursday 5 November 2026}. You have four
books (Modules 5 to 8). Here is a simple plan.

\begin{table}
  \centering
  \begin{tabular}{@{}L{30mm}L{\dimexpr\linewidth-34mm\relax}@{}}
    \toprule
    \colhead{When} & \colhead{What to do}\\
    \midrule
    Days 1 to 10 & Work through the four books, about two or three days for
      each book. For every chapter, follow steps 1 to 5 below.\\
    Days 11 to 13 & Do the Practice set of each book like a real exam: time
      yourself and use only the formula sheet. Mark your answers with the
      marking criteria. Redo every question you got wrong.\\
    The night before & Read the Last-minute summary of each book and look at
      the formula sheet once more. Then go to bed early.\\
    \bottomrule
  \end{tabular}
\end{table}

For each chapter:
\begin{enumerate}
  \item Read the \textbf{big idea} and the \textbf{key words} (about 10 minutes).
  \item Read the explanation and the \textbf{formula cards}. Say each formula
        ``in words''.
  \item \textbf{Cover the solution} of each worked example and try it yourself.
        Then compare your work with the book, step by step.
  \item Do the \textbf{Quick check} questions. Mark them with the answers at the
        back.
  \item Write every mistake you made in a ``mistakes list''. Read this list
        again before the exam.
\end{enumerate}

\section*{How to use the official formula sheet in the exam}

In the exam you get the NESA data sheet and formulae sheet. It gives many
formulas and constants, but not all of them.

\begin{enumerate}
  \item \textbf{Before the exam}, learn where each formula is on the sheet. The
        sheet has four parts: motion, forces and gravity; waves and
        thermodynamics; electricity and magnetism; quantum, special relativity
        and nuclear. Then you can find a formula quickly.
  \item \textbf{Learn the formulas that are not on the sheet.} Their formula
        cards have the red badge and show you how to derive them.
  \item \textbf{Check what each symbol means} in the question. The same letter
        can mean different things: for example, $T$ is a period in one formula
        and a temperature in another.
  \item \textbf{Change every quantity to SI units first}: km to m, nm to m,
        \si{km.h^{-1}} to \si{m.s^{-1}}, eV to J.
  \item \textbf{Use the values on the data sheet}, for example
        $g = \qty{9.8}{m.s^{-2}}$ and $c = \qty{3.00e8}{m.s^{-1}}$.
  \item \textbf{Show your working} in the same order as the worked examples:
        formula, substitution, answer with a unit.
\end{enumerate}

% Numbers in the mini worked example: R = 12 / 0.50 = 24 ohm (checked in Python).
\unnumberedchapter[\faInfoCircle]{How to use this guide}

This guide explains Module~\modulenumber\ (\moduletitle) from the very beginning.
It uses simple English, short sentences and many examples. Every calculation is
shown step by step.

The chapters follow the inquiry questions of the NSW Physics Stage 6 Syllabus
(2017). At the back of the book you will find a formula summary, a glossary
(with Chinese translations), a practice set, full answers and a last-minute
summary. In the PDF, you can click any entry in the contents, any
cross-reference and any question number.

\section*{What the coloured boxes mean}

Each kind of box has its own colour and icon. A colour always means the same
thing, in every book of this series.

\begin{sidebyside}[0.5]
  \begin{bigidea}[before skip=0pt, after skip=0pt, equal height group=legend1, fontupper=\normalsize]
    The main point of a topic in two or three sentences. Read these first.
  \end{bigidea}
\nextside
  \begin{keywords}[before skip=0pt, after skip=0pt, equal height group=legend1]
    \kw{Term}{what the word means, in plain words.}
    New words are also printed in \term{bold} the first time they appear.
  \end{keywords}
\end{sidebyside}

\begin{sidebyside}[0.5]
  \begin{mistakes}[before skip=0pt, after skip=0pt, equal height group=legend2]
    Errors that students often make, and how to avoid them:
    \xmark\ wrong \enspace \cmark\ right.
  \end{mistakes}
\nextside
  \begin{examtip}[before skip=0pt, after skip=0pt, equal height group=legend2]
    How the HSC asks about this topic, and what you must write to get full
    marks.
  \end{examtip}
\end{sidebyside}

\begin{sidebyside}[0.5]
  \begin{toolkit}[before skip=0pt, after skip=0pt, equal height group=legend3]{}
    A short maths skill that you need for the next part, for example
    rearranging a formula.
  \end{toolkit}
\nextside
  \begin{refresher}[before skip=0pt, after skip=0pt, equal height group=legend3]{}
    A short reminder of a Year~11 idea that you need, for example Newton's
    laws.
  \end{refresher}
\end{sidebyside}

\begin{sidebyside}[0.5]
  \begin{tcolorbox}[hsc minor=hscqc, unbreakable, before skip=0pt, after skip=0pt,
      equal height group=legend4, title={\hscboxtitle{\faQuestionCircle}{Quick check}{}}]
    Short questions to test yourself, numbered like \textcolor{hscqc}{\bfseries Q2.3}
    (Chapter~2, question~3). The answers at the back use the same numbers.
  \end{tcolorbox}
\nextside
  \begin{experiment}[unbreakable, before skip=0pt, after skip=0pt, equal height group=legend4]{}
    A practical or historical experiment: aim, apparatus, method, results,
    conclusion and why it matters.
  \end{experiment}
\end{sidebyside}

\begin{formulacard}[unbreakable, before skip=6pt]{Name of the formula}
  \cardformula{V = IR}
  Every formula has a card: the formula, the meaning and SI unit of each
  symbol, when to use it, and a common trap. The badge under the formula tells
  you about the official HSC formula sheet:
  \begin{itemize}
    \item \hscbadgeonsheet\ \ it is printed on the sheet, so you are given it
          in the exam;
    \item \hscbadgenotonsheet\ \ you must learn it, or know how to get it;
    \item \hscbadgecombine{two formulas}\ \ it is two formulas from the sheet
          joined together.
  \end{itemize}
\end{formulacard}

\begin{workedexample}[nonumber, unbreakable]{How every worked example is set out}
  A \qty{12}{V} battery drives a current of \qty{0.50}{A} through a resistor.
  Find the resistance. (Cover everything below the \emph{Solution} line with a
  sheet of paper and try it first.)
  \given{$V = \qty{12}{V}$, \ $I = \qty{0.50}{A}$}
  \find{the resistance $R$}
  \formula{$V = IR$, \ so $R = \dfrac{V}{I}$}
  \substitute{$R = \dfrac{\qty{12}{V}}{\qty{0.50}{A}}$ \ (numbers with their units)}
  \answer{\final{$R = \qty{24}{\ohm}$} \ (with a unit and sensible significant
    figures)}
  \checkstep{is the answer sensible? Units: $\unit{V} \div \unit{A} = \unit{\ohm}$ \cmark}
\end{workedexample}

Each chapter starts with its \textbf{inquiry question} from the syllabus, a
list of \textbf{goals} and the \textbf{syllabus dot points} it covers. It ends
with a short \textbf{summary}.

\section*{How to study with this guide in the last two weeks}

The HSC Physics exam is on \textbf{Thursday 5 November 2026}. You have four
books (Modules 5 to 8). Here is a simple plan.

\begin{table}
  \centering
  \begin{tabular}{@{}L{30mm}L{\dimexpr\linewidth-34mm\relax}@{}}
    \toprule
    \colhead{When} & \colhead{What to do}\\
    \midrule
    Days 1 to 10 & Work through the four books, about two or three days for
      each book. For every chapter, follow steps 1 to 5 below.\\
    Days 11 to 13 & Do the Practice set of each book like a real exam: time
      yourself and use only the formula sheet. Mark your answers with the
      marking criteria. Redo every question you got wrong.\\
    The night before & Read the Last-minute summary of each book and look at
      the formula sheet once more. Then go to bed early.\\
    \bottomrule
  \end{tabular}
\end{table}

For each chapter:
\begin{enumerate}
  \item Read the \textbf{big idea} and the \textbf{key words} (about 10 minutes).
  \item Read the explanation and the \textbf{formula cards}. Say each formula
        ``in words''.
  \item \textbf{Cover the solution} of each worked example and try it yourself.
        Then compare your work with the book, step by step.
  \item Do the \textbf{Quick check} questions. Mark them with the answers at the
        back.
  \item Write every mistake you made in a ``mistakes list''. Read this list
        again before the exam.
\end{enumerate}

\section*{How to use the official formula sheet in the exam}

In the exam you get the NESA data sheet and formulae sheet. It gives many
formulas and constants, but not all of them.

\begin{enumerate}
  \item \textbf{Before the exam}, learn where each formula is on the sheet. The
        sheet has four parts: motion, forces and gravity; waves and
        thermodynamics; electricity and magnetism; quantum, special relativity
        and nuclear. Then you can find a formula quickly.
  \item \textbf{Learn the formulas that are not on the sheet.} Their formula
        cards have the red badge and show you how to derive them.
  \item \textbf{Check what each symbol means} in the question. The same letter
        can mean different things: for example, $T$ is a period in one formula
        and a temperature in another.
  \item \textbf{Change every quantity to SI units first}: km to m, nm to m,
        \si{km.h^{-1}} to \si{m.s^{-1}}, eV to J.
  \item \textbf{Use the values on the data sheet}, for example
        $g = \qty{9.8}{m.s^{-2}}$ and $c = \qty{3.00e8}{m.s^{-1}}$.
  \item \textbf{Show your working} in the same order as the worked examples:
        formula, substitution, answer with a unit.
\end{enumerate}
}{}
\begingroup
  \contentsmargin[0.65em]{1.9em}%
  \tableofcontents
\endgroup

\unnumberedchapter[\faMap]{Module overview}

% ---- helpers for the overview only (prefix ov)
% sheet marks: on the sheet / learn it / combine (same colours and icons as \frow)
\newcommand\ovYes{\textcolor{hscformula}{\faCheck}}
\newcommand\ovNo{\textcolor{hscmistake}{\faTimes}}
\newcommand\ovComb{\textcolor{hsctiptext}{\faLink}}
% one row of the formula grid:  \ovrow{mark}{formula}{what it is for}
\newcommand\ovrow[3]{#1 & #2 & {\footnotesize #3}\\\hline}
% chapter heading row of the formula grid
\newcommand\ovch[2]{\multicolumn{3}{@{}l@{}}{\rule{0pt}{3.4ex}\cref{#1}:
  {\bfseries\color{hscbrand}#2}}\\\hline}
\newenvironment{ovgrid}{%
  \small\renewcommand\arraystretch{1.32}\arrayrulecolor{hscrule}%
  \begin{tabular}[t]{@{}C{4.5mm}@{\hspace{1.5mm}}L{40mm}@{\hspace{1.5mm}}L{\dimexpr\linewidth-47.5mm\relax}@{}}}%
  {\end{tabular}}

\section*{What this module is about}

Module~6 is about how electricity and magnetism work together. Two ideas run through the
whole module.
\begin{itemize}
  \item A magnetic field pushes on a \emph{moving} charge. So a wire that carries a current
        feels a force in a magnetic field. This is the \emph{motor effect}.
  \item A \emph{changing} magnetic field makes a voltage (an emf) in a wire or a coil. This
        is \emph{electromagnetic induction}.
\end{itemize}
These two ideas explain motors, generators, transformers, the power grid and magnetic
brakes. In Year~11 you met charge, current, voltage and the magnetic field of a current.
This module puts them to work.

\section*{How the parts fit together}

\Cref{ov:fig:map} shows how the five chapters build on each other.

\begin{figure}
  \centering
  \begin{tikzpicture}[x=1cm, y=1cm]
    \node[concept=104mm, fill=hsctool!8, draw=hsctool!60] (y11) at (0,0)
      {\textbf{Year 11 tools:} charge, current, voltage, force, energy;
       a current makes a magnetic field};
    \node[concept=52mm] (c1) at (-4.0,-1.85)
      {\textbf{Ch 1 Charged particles in fields}\\ \mbox{$F = qE$} and \mbox{$F = qvB\sin\theta$}};
    \node[concept=52mm] (c2) at (4.0,-1.85)
      {\textbf{Ch 2 The motor effect}\\ \mbox{$F = lIB\sin\theta$}; parallel wires};
    \node[concept=52mm] (c3) at (-4.0,-4.0)
      {\textbf{Ch 3 Electromagnetic induction}\\ flux; Faraday's and Lenz's laws};
    \node[concept=52mm] (c5) at (4.0,-4.0)
      {\textbf{Ch 5 Motors, generators, brakes}\\ torque; back emf; eddy currents};
    \node[concept=62mm] (c4) at (-4.0,-6.15)
      {\textbf{Ch 4 Transformers and power lines}\\ step up, step down; \mbox{$I^2R$} loss};
    \draw[link] (c1.north |- y11.south) -- node[linklabel, left=5pt] {forces on charges} (c1.north);
    \draw[link] (c2.north |- y11.south) -- node[linklabel, right=5pt] {field of a current} (c2.north);
    \draw[link] (c1.east) -- node[linklabel, above=1pt, align=center] {current = many\\moving charges} (c2.west);
    \draw[link] (c1.south) -- node[linklabel, right=5pt, align=left] {charges in a\\moving wire} (c3.north);
    \draw[link] (c2.south) -- node[linklabel, left=5pt, align=right] {forces on a coil\\give a torque} (c5.north);
    \draw[link] (c3.east) -- node[linklabel, above=1pt, align=center] {changing flux in\\a turning coil} (c5.west);
    \draw[link] (c3.south) -- node[linklabel, right=5pt, align=left] {two coils share\\a changing flux} (c4.north);
    \draw[link] (c5.south) |- node[linklabel, pos=0.75, above=1pt] {generators make the AC} (c4.east);
  \end{tikzpicture}
  \caption{How the ideas in this module are linked. Each arrow means ``this idea is used in
    the next one'' (``Ch'' means chapter).}\label{ov:fig:map}
\end{figure}

\section*{The inquiry questions}

The syllabus asks four inquiry questions (big questions). \Cref{ov:tab:iq} shows which
chapters answer each one.

\begin{table}
  \caption{The inquiry questions of Module~6 and the chapters that answer them.}
  \label{ov:tab:iq}
  \centering
  \small
  \begin{tabular}{@{}L{10mm}L{97mm}L{33mm}@{}}
    \toprule
    & \colhead{Inquiry question (from the syllabus)} & \colhead{Answered in}\\
    \midrule
    \textbf{IQ1} & \emph{What happens to stationary and moving charged particles when they
      interact with an electric or magnetic field?} & \cref{ch1:chap}\\
    \textbf{IQ2} & \emph{Under what circumstances is a force produced on a current-carrying
      conductor in a magnetic field?} & \cref{ch2:chap}\\
    \textbf{IQ3} & \emph{How are electric and magnetic fields related?}
      & \cref{ch3:chap,ch4:chap}\\
    \textbf{IQ4} & \emph{How has knowledge about the motor effect been applied to
      technological advances?} & \cref{ch5:chap}\\
    \bottomrule
  \end{tabular}
\end{table}

\section*{Key experiments}

The syllabus asks you to carry out or study these investigations (\cref{ov:tab:experiments}).
Each has a Key experiment box: aim, method, results, conclusion, and validity, reliability
and accuracy.

\begin{table}
  \caption{The key experiments of Module~6. ``Required'' marks the practical investigation
    that the syllabus names.}\label{ov:tab:experiments}
  \centering
  \small
  \begin{tabular}{@{}L{46mm}L{56mm}L{17mm}L{17mm}@{}}
    \toprule
    \colhead{Experiment} & \colhead{What it shows} & \colhead{Dot point} & \colhead{Where}\\
    \midrule
    Paths of charges between plates (simulation)
      & the path is a parabola; the deflection grows in step with $V$
      & M6.IQ1.01, IQ1.02 & \cref{ch1:sec:experiments}\\
    Speed gained from an accelerating voltage (simulation)
      & $v^2$ grows in step with $V$, as $qV = \frac{1}{2}mv^2$ predicts
      & M6.IQ1.01 & \cref{ch1:sec:experiments}\\
    The force on a wire (current balance)
      & $F$ grows in step with $I$; largest at \ang{90} to $B$, zero when parallel
      & M6.IQ2.01 & \cref{ch2:sec:experiments}\\
    Two parallel wires (\textbf{required})
      & like currents attract, unlike repel; testing $F/l = (\mu_0/2\pi)(I_1 I_2/r)$
      & M6.IQ2.02 & \cref{ch2:sec:experiments}\\
    Faraday's induction experiments (1831)
      & an emf is induced only while the flux changes
      & M6.IQ3.02 & \cref{ch3:sec:experiments}\\
    Falling magnets and swinging plates
      & eddy currents oppose the motion (Lenz); energy becomes heat
      & M6.IQ3.02, IQ4.03 & \cref{ch3:sec:experiments}\\
    Testing the transformer equations
      & $V_\mathrm{s}$ grows in step with $N_\mathrm{s}$; real transformers lose some energy
      & M6.IQ3.03, IQ3.04 & \cref{ch4:sec:experiments}\\
    Investigating a simple DC motor
      & the job of each part; back emf lowers the current
      & M6.IQ4.01 & \cref{ch5:sec:experiments}\\
    \bottomrule
  \end{tabular}
\end{table}

\section*{All the formulas at a glance}

Every formula card in this book. \ovYes\ = on the HSC formula sheet; \ovNo\ = not on the
sheet: learn it or derive it; \ovComb\ = join formulas from the sheet.
\hyperref[fs:chap]{\emph{Formula sheet explained}} (page~\pageref{fs:chap}) says when to use
each one.

\par\smallskip\noindent
\begin{minipage}[t]{0.488\linewidth}
\begin{ovgrid}
  \ovch{ch1:chap}{Charged particles in fields}
  \ovrow{\ovYes}{$E = \frac{V}{d}$}{field between parallel plates}
  \ovrow{\ovYes}{$\vec{F} = q\vec{E}$}{force on a charge in an electric field}
  \ovrow{\ovNo}{$qE = mg$, \ $q = ne$}{charged drop held still between plates}
  \ovrow{\ovYes}{$W = qV$, \ $W = qEd$}{work done on a charge by the field}
  \ovrow{\ovComb}{$\frac{1}{2}mv^2 = qV$\newline $v = \sqrt{\frac{2qV}{m}}$}
    {speed gained from rest}
  \ovrow{\ovComb}{$t = \frac{L}{u}$\newline $y = \frac{1}{2}at^2 = \frac{qEL^2}{2mu^2}$}
    {sideways deflection between plates}
  \ovrow{\ovYes}{$F = qv_\perp B = qvB\sin\theta$}{magnetic force on a moving charge}
  \ovrow{\ovNo}{$r = \frac{mv}{qB}$, \ $T = \frac{2\pi m}{qB}$}
    {radius and period of the circle in a magnetic field}
  \ovch{ch2:chap}{The motor effect}
  \ovrow{\ovYes}{$B = \frac{\mu_0 I}{2\pi r}$}{field near a long straight wire}
  \ovrow{\ovYes}{$B = \frac{\mu_0 NI}{L}$}{field inside a solenoid}
  \ovrow{\ovYes}{$F = lI_\perp B = lIB\sin\theta$}{force on a current-carrying wire}
  \ovrow{\ovYes}{$\frac{F}{l} = \frac{\mu_0}{2\pi}\frac{I_1 I_2}{r}$}
    {force between two parallel wires}
\end{ovgrid}
\end{minipage}\hfill
\begin{minipage}[t]{0.488\linewidth}
\begin{ovgrid}
  \ovch{ch3:chap}{Electromagnetic induction}
  \ovrow{\ovYes}{$\Phi = B_\parallel A = BA\cos\theta$}{magnetic flux through a loop}
  \ovrow{\ovYes}{$\varepsilon = -N\frac{\Delta\Phi}{\Delta t}$}
    {Faraday's law: the induced emf}
  \ovrow{\ovNo}{$\varepsilon = lvB$}{emf of a rod moving across a field}
  \ovrow{\ovComb}{$\lvert\varepsilon_\mathrm{s}\rvert = N_\mathrm{s}\frac{\mu_0 N_\mathrm{p} A}{L}
    \frac{\lvert\Delta I_\mathrm{p}\rvert}{\Delta t}$}{emf in a coil wound on a solenoid}
  \ovch{ch4:chap}{Transformers and power lines}
  \ovrow{\ovYes}{$\frac{V_\mathrm{p}}{V_\mathrm{s}} = \frac{N_\mathrm{p}}{N_\mathrm{s}}$}
    {transformer (turns) equation}
  \ovrow{\ovYes}{$P = VI$}{electrical power}
  \ovrow{\ovYes}{$V_\mathrm{p} I_\mathrm{p} = V_\mathrm{s} I_\mathrm{s}$}
    {ideal transformer: power in = power out}
  \ovrow{\ovNo}{$\eta = \frac{P_\mathrm{out}}{P_\mathrm{in}} \times 100\%$\newline
    $\phantom{\eta} = \frac{V_\mathrm{s}I_\mathrm{s}}{V_\mathrm{p}I_\mathrm{p}} \times 100\%$}
    {efficiency of a real transformer}
  \ovrow{\ovYes}{$P = \frac{\Delta E}{\Delta t}$}{energy used or wasted in a time}
  \ovrow{\ovNo}{$P_\mathrm{loss} = I^2R$}{power lost as heat in a line}
  \ovch{ch5:chap}{Motors, generators, brakes}
  \ovrow{\ovYes}{$\tau = nIA_\perp B = nIAB\sin\theta$}{torque on a coil}
  \ovrow{\ovNo}{$I = \frac{V - \varepsilon_\mathrm{back}}{R}$}
    {current in a running motor (back emf)}
  \ovrow{\ovYes}{$\varepsilon = -N\frac{\Delta\Phi}{\Delta t}$}
    {average emf of a turning generator coil}
  \ovrow{\ovComb}{$Q = \frac{1}{2}mu^2 - \frac{1}{2}mv^2$}{heat made by a magnetic brake}
\end{ovgrid}
\end{minipage}
\par\medskip
